\\documentclass{article}
\\usepackage{amsmath}
\\usepackage{amsfonts}
\\usepackage{amssymb}
\\begin{document}

\\section*{Warp Bubble Parameter Constraints Analysis}

\\subsection*{Model}
We analyze a simplified warp bubble with Gaussian profile:
\\[
  f(r,t) = 1 + A e^{-\\frac{(r-R)^2}{\\sigma^2}} e^{-\\frac{t^2}{\\sigma^2}}
\\]

The corresponding energy density is:
\\[
  T^{00}(r,t) = -\\frac{A^2}{8\\pi r^2} e^{-\\frac{2(r-R)^2}{\\sigma^2}} e^{-\\frac{2t^2}{\\sigma^2}}
\\]

\\subsection*{Total Negative Energy}
The total negative energy within the warp bubble region is:
\\[
  E_{\\rm neg}(R,\\sigma,A) = 2 \pi \left(- \frac{\sqrt{2} A^{2} \sigma \operatorname{erf}{\left(\frac{\sqrt{2} R}{\sigma} \right)}}{16 \sqrt{\pi}} - \frac{\sqrt{2} A^{2} \sigma \operatorname{erf}{\left(\frac{2 \sqrt{2} R}{\sigma} \right)}}{16 \sqrt{\pi}}\right)
\\]

\\subsection*{Extremum Conditions}
\\[
  \\frac{\\partial E_{\\rm neg}}{\\partial R} = 2 \pi \left(- \frac{A^{2} e^{- \frac{2 R^{2}}{\sigma^{2}}}}{4 \pi} - \frac{A^{2} e^{- \frac{8 R^{2}}{\sigma^{2}}}}{2 \pi}\right) = 0
\\]
\\[
  \\Longrightarrow R = \text{No analytical solutions found}
\\]

\\[
  \\frac{\\partial E_{\\rm neg}}{\\partial \\sigma} = 2 \pi \left(\frac{A^{2} R e^{- \frac{2 R^{2}}{\sigma^{2}}}}{4 \pi \sigma} + \frac{A^{2} R e^{- \frac{8 R^{2}}{\sigma^{2}}}}{2 \pi \sigma} - \frac{\sqrt{2} A^{2} \operatorname{erf}{\left(\frac{\sqrt{2} R}{\sigma} \right)}}{16 \sqrt{\pi}} - \frac{\sqrt{2} A^{2} \operatorname{erf}{\left(\frac{2 \sqrt{2} R}{\sigma} \right)}}{16 \sqrt{\pi}}\right) = 0
\\]
\\[
  \\Longrightarrow \\sigma = \text{No analytical solutions found}
\\]

\\subsection*{Nondimensional Analysis}

Defining the dimensionless parameter $x = R/\\sigma$, the extremum condition 
$\\partial E_{\\rm neg}/\\partial R = 0$ reduces to the transcendental equation:
\\[
  F(x) = e^{-2x^2} + \frac{1}{2}e^{-8x^2} = 0
\\]

Since both exponential terms are strictly positive for all $x > 0$, we have $F(x) > 0$ 
for all real $x$. This indicates that \\textbf{no critical points exist} for this 
simplified model.

\\subsection*{Asymptotic Analysis}

\\textbf{Thin-wall limit} ($\\sigma \\ll R$, i.e., $x \\gg 1$):
\\[
  \\text{erf}(x) \\approx 1 - \\frac{e^{-x^2}}{\\sqrt{\\pi} x}
\\]

\\textbf{Thick-wall limit} ($\\sigma \\gg R$, i.e., $x \\ll 1$):
\\[
  \\text{erf}(x) \\approx \\frac{2x}{\\sqrt{\\pi}}\\left(1 - \\frac{x^2}{3} + \\cdots\\right)
\\]

\\subsection*{Physical Constraints}
For exotic matter generation and finite energy:
\\[
  A > 0\ \text{ for exotic matter generation}
\\]

\\subsection*{Physical Interpretation}

The absence of critical points in this simplified model suggests:
\\begin{itemize}
  \\item The energy $E_{\\rm neg}$ decreases monotonically with increasing $R$ or $\\sigma$
  \\item No natural ``optimal'' bubble size exists
  \\item Physical constraints (finite resources, stability requirements) must determine parameters
  \\item More sophisticated models may include additional terms that create energy minima
\\end{itemize}

\\subsection*{Numerical Approaches}

For practical parameter optimization, consider:
\\begin{enumerate}
  \\item \\textbf{Parameter scanning}: Grid search over $(R,\\sigma)$ space
  \\item \\textbf{Constrained optimization}: Minimize $|E_{\\rm neg}|$ subject to physical bounds
  \\item \\textbf{Multi-objective optimization}: Balance energy efficiency with stability constraints
\\end{enumerate}

\\subsection*{Notes}
The extremum conditions for this model lead to transcendental equations involving error functions
that cannot be solved analytically in general. The nondimensional analysis reveals that
the simplified Gaussian model has no critical points, indicating the need for more 
sophisticated warp bubble models that include additional physical effects.

\\end{document}
